\needspace{8\baselineskip}
\subsection*{Constructing and interpreting singular values}

\begin{problem}\label{ex:10:svd-left}
For a real \(m\times n\) matrix \(A\) of rank \(r>0\), construct matrices
\(S\in\R^{m\times r}\), \(T\in\R^{n\times r}\), with orthonormal columns,
and a positive diagonal \(r\times r\) matrix \(\Lambda\) such that
\(A=S\Lambda^{1/2}T^\top\). Start with the eigenvectors of \(AA^\top\).
\end{problem}
\begin{solution}
Choose orthonormal eigenvectors \(s_1,\ldots,s_r\) for the positive
eigenvalues of \(AA^\top\), and let \(S\) have these columns and
\(\Lambda\) these eigenvalues. Define \(T=A^\top S\Lambda^{-1/2}\).
Then
\[
T^\top T=\Lambda^{-1/2}S^\top AA^\top S\Lambda^{-1/2}=I_r.
\]
The remaining orthonormal eigenvectors form \(S_0\); for each such vector
\(z\), \(\norm{A^\top z}^2=z^\top AA^\top z=0\), so \(S_0^\top A=0\).
Consequently
\[
A=(SS^\top+S_0S_0^\top)A=SS^\top A=S\Lambda^{1/2}T^\top.
\]
\end{solution}

\begin{problem}\label{ex:10:linked-vectors}
In an SVD, the equations \(AA^\top S=S\Lambda\) and
\(A^\top AT=T\Lambda\) hold. Explain why solving these two equations
independently, even with orthonormal columns, need not give
\(A=S\Lambda^{1/2}T^\top\). Give the correct procedure.
\end{problem}
\begin{solution}
The equations do not link the two choices of signs or repeated-eigenspace
bases. For \(A=(1)\), \(S=(1)\), \(T=(-1)\), \(\Lambda=(1)\), both
eigenvalue equations hold but the proposed product is \(-1\).
After choosing \(S,\Lambda\) from \(AA^\top\), set
\(T=A^\top S\Lambda^{-1/2}\), as in the preceding proof.
Alternatively choose \(T\) first and set \(S=AT\Lambda^{-1/2}\).
For a repeated positive singular value the same orthogonal change must
be applied to both corresponding lists.
\end{solution}

\begin{problem}\label{ex:10:polar}
Prove the two full-rank polar decompositions:
(a) if \(A\in\R^{m\times n}\) has full column rank, then \(A=PH\) with
\(P^\top P=I_n\) and \(H\) positive definite;
(b) if \(A\) has full row rank, then \(A=HP^\top\) with \(P^\top P=I_m\)
and \(H\) positive definite.
\end{problem}
\begin{solution}
(a) Set \(H=(A^\top A)^{1/2}\), positive definite by full column rank,
and \(P=AH^{-1}\). Then \(PH=A\) and
\(P^\top P=H^{-1}A^\top AH^{-1}=I_n\).
(b) Apply (a) to \(A^\top\): write \(A^\top=PH\), where now
\(H=(AA^\top)^{1/2}\) and \(P^\top P=I_m\). Transposition gives
\(A=HP^\top\), since \(H=H^\top\).
\end{solution}

\begin{problem}\label{ex:10:singular-order}
Let \(A,B\) be symmetric, with \(A\), \(B\), and \(A-B\) positive
semidefinite. If \(A\) is singular, prove that \(B\) and \(A-B\) are singular.
\end{problem}
\begin{solution}
Choose nonzero \(z\in\ker A\). Then
\(0=z^\top Az=z^\top Bz+z^\top(A-B)z\), a sum of nonnegative terms,
so both vanish. For a positive semidefinite \(C\),
\(z^\top Cz=\norm{C^{1/2}z}^2=0\) implies \(Cz=0\).
Apply this to \(B,A-B\); both have the nonzero kernel vector \(z\).
\end{solution}

\begin{problem}\label{ex:10:sqrt-dependence}
For a positive semidefinite \(A\), prove that \(A^{1/2}x,A^{1/2}y\)
are dependent if and only if \(Ax,Ay\) are dependent.
\end{problem}
\begin{solution}
An orthonormal eigenbasis shows \(\ker A=\ker A^{1/2}\), since a
nonnegative eigenvalue and its square root vanish together.
For any \((a,b)\ne(0,0)\),
\[
aAx+bAy=0\iff ax+by\in\ker A
\iff ax+by\in\ker A^{1/2}\iff aA^{1/2}x+bA^{1/2}y=0.
\]
Thus a nontrivial relation exists for one pair exactly when it exists for the other.
\end{solution}

\begin{problem}\label{ex:10:numerical-values}
Find all eigenvalues, including nonreal ones, and singular values of
\[
\text{(a) }\begin{pmatrix}0&1\\0&0\end{pmatrix},\quad
\text{(b) }\begin{pmatrix}1&1\\0&0\end{pmatrix},\quad
\text{(c) }\begin{pmatrix}1&1\\1&1\end{pmatrix},\quad
\text{(d) }\begin{pmatrix}0&1\\-1&0\end{pmatrix}.
\]
\end{problem}
\begin{solution}
(a) The characteristic polynomial is \(t^2\), so both eigenvalues are
zero; \(A^\top A=\operatorname{diag}(0,1)\), giving singular values \(1,0\).
(b) The triangular matrix has eigenvalues \(1,0\);
\(A^\top A=\begin{pmatrix}1&1\\1&1\end{pmatrix}\) has eigenvectors
\((1,1)^\top,(1,-1)^\top\) with eigenvalues \(2,0\). The singular values
are \(\sqrt2,0\).
(c) Those same directions give eigenvalues \(2,0\) for \(A\), and
\(A^\top A=2A\), so its eigenvalues are \(4,0\). The singular values are \(2,0\).
(d) The characteristic polynomial factors as \(t^2+1=(t-i)(t+i)\),
so the eigenvalues are \(i,-i\). For example,
\(A(1,i)^\top=(i,-1)^\top=i(1,i)^\top\); conjugating gives the other
eigenvector. Since \(A^\top A=I_2\), both singular values are \(1\).
Thus the nonreal eigenvalues of this real rotation coexist with real,
nonnegative singular values.
\end{solution}

\begin{problem}\label{ex:10:transpose-numerical}
Let
\[
A=\begin{pmatrix}1&1&0\\1&1&0\\1&-1&-1\end{pmatrix}.
\]
Find the eigenvalues and singular values of \(A\) and \(A^\top\).
For \(A^\top\), list the eigenvalues, their absolute values, and singular
values in decreasing order.
\end{problem}
\begin{solution}
Block triangularity gives the characteristic polynomial
\(t(t-2)(t+1)\). Transposition preserves this polynomial.
Furthermore,
\[
A^\top A=\begin{pmatrix}3&1&-1\\1&3&1\\-1&1&1\end{pmatrix}.
\]
The mutually orthogonal vectors \((1,1,0)^\top,(1,-1,-1)^\top,
(1,-1,2)^\top\) give eigenvalues \(4,3,0\), respectively, by direct
multiplication. Thus both \(A\) and \(A^\top\) have singular values
\(2,\sqrt3,0\). The three decreasing lists are
\(2,0,-1\); \(2,1,0\); and \(2,\sqrt3,0\).
\end{solution}

\begin{problem}\label{ex:10:spectral-versus-singular}
For \(A=\begin{pmatrix}1&2\\2&1\end{pmatrix}\), compute the eigenvalues,
singular values, a spectral decomposition, and the modulus
\(|A|=(A^\top A)^{1/2}\).
\end{problem}
\begin{solution}
The orthonormal vectors \(q_1=(1,1)^\top/\sqrt2\),
\(q_2=(1,-1)^\top/\sqrt2\) have eigenvalues \(3,-1\), so
\(A=3q_1q_1^\top-q_2q_2^\top\).
Squaring gives \(A^\top A=9q_1q_1^\top+q_2q_2^\top\), with singular
values \(3,1\). Its positive semidefinite square root is
\[
|A|=3q_1q_1^\top+q_2q_2^\top=\begin{pmatrix}2&1\\1&2\end{pmatrix}.
\]
The negative eigenvalue of \(A\) becomes a positive singular value.
\end{solution}

\begin{problem}\label{ex:10:full-svd-construction}
Find a Schur triangular form (an upper triangular matrix obtained by
orthogonal similarity) and a full SVD for
\[
A=\begin{pmatrix}1&0&0&1\\0&2&3&0\\0&0&2&0\\1&0&0&1\end{pmatrix}.
\]
\end{problem}
\begin{solution}
The orthonormal basis \((e_1+e_4)/\sqrt2,(e_1-e_4)/\sqrt2,e_2,e_3\)
gives the upper triangular matrix
\(\operatorname{diag}\left(2,0,\begin{pmatrix}2&3\\0&2\end{pmatrix}\right)\).
Indeed the first two basis vectors have eigenvalues \(2,0\), while
\(Ae_2=2e_2\) and \(Ae_3=3e_2+2e_3\).
For the remaining \(2\times2\) block, its Gram matrix is
\(\begin{pmatrix}4&6\\6&13\end{pmatrix}\); multiplying \((1,2)^\top\)
and \((-2,1)^\top\) gives eigenvalues \(16,1\).
Thus, in decreasing singular-value order, choose
\[
\begin{array}{ll}
v_1=(0,1,2,0)^\top/\sqrt5,&u_1=(0,2,1,0)^\top/\sqrt5,\\
v_2=(1,0,0,1)^\top/\sqrt2,&u_2=v_2,\\
v_3=(0,-2,1,0)^\top/\sqrt5,&u_3=(0,-1,2,0)^\top/\sqrt5,\\
v_4=(1,0,0,-1)^\top/\sqrt2,&u_4=v_4.
\end{array}
\]
Both lists are orthonormal: their two coordinate planes are perpendicular,
and the pairs within each plane have zero dot product and unit norm.
Direct multiplication gives \(Av_1=4u_1\), \(Av_2=2u_2\),
\(Av_3=u_3\), \(Av_4=0\). Hence, with the displayed columns,
\(A=U\operatorname{diag}(4,2,1,0)V^\top\), the required full SVD.
\end{solution}

\needspace{8\baselineskip}
\subsection*{Moore--Penrose inverses and their identities}

\begin{problem}\label{ex:10:mp-basic}
Find the Moore--Penrose inverse of (a) a nonsingular square matrix;
(b) a scalar; (c) a diagonal matrix; (d) an \(m\times n\) zero matrix.
\end{problem}
\begin{solution}
(a) \(A^{-1}\): both products are \(I\), and the two triple products
are \(A\) and \(A^{-1}\).
(b) For scalar \(a\), use \(a^+=1/a\) when \(a\ne0\), and zero when
\(a=0\). Each scalar then satisfies the four equations.
(c) For \(D=\operatorname{diag}(d_i)\), take
\(D^+=\operatorname{diag}(d_i^+)\). The triple equations hold coordinatewise;
\(DD^+=D^+D\) is diagonal with entries zero or one, hence symmetric.
(d) The inverse is \(0_{n\times m}\). All four equations are zero
equalities of the appropriate sizes.
\end{solution}

\begin{problem}\label{ex:10:mp-existence}
Prove existence of the Moore--Penrose inverse by checking all four equations.
\end{problem}
\begin{solution}
If \(A=U_rDV_r^\top\), put \(X=V_rD^{-1}U_r^\top\).
Then \(AX=U_rU_r^\top\) and \(XA=V_rV_r^\top\), both symmetric.
Moreover
\(AXA=U_rU_r^\top U_rDV_r^\top=A\) and
\(XAX=V_rV_r^\top V_rD^{-1}U_r^\top=X\).
When \(A=0\), the zero matrix of reversed size satisfies every equation.
\end{solution}

\begin{problem}\label{ex:10:mp-uniqueness}
Prove uniqueness using only the four Moore--Penrose equations.
\end{problem}
\begin{solution}
Let \(B,C\) both satisfy the equations for \(A\).
Symmetry and \(ACA=A\) give
\[
AB=B^\top A^\top=B^\top(ACA)^\top
=(AB)^\top(AC)^\top=ABAC=AC,
\]
where \(ABA=A\) gives the last equality. Likewise
\[
BA=A^\top B^\top=(ACA)^\top B^\top
=(CA)^\top(BA)^\top=CABA=CA.
\]
Now \(B=BAB=(BA)B=(CA)B=C(AB)=C(AC)=CAC=C\).
\end{solution}

\begin{problem}\label{ex:10:mp-transpose}
Prove \((A^+)^+=A\) and \((A^\top)^+=(A^+)^\top\).
\end{problem}
\begin{solution}
The four equations are unchanged when \(A,A^+\) are interchanged:
the triple equations exchange roles, and the two symmetry requirements
exchange roles. Uniqueness therefore gives \((A^+)^+=A\).
Transposing \(AA^+A=A\) and \(A^+AA^+=A^+\) gives the two triple
equations for \(A^\top,(A^+)^\top\).
Their products are \((A^+A)^\top\) and \((AA^+)^\top\), which are
symmetric. Uniqueness proves the second identity.
\end{solution}

\begin{problem}\label{ex:10:mp-projectors}
For \(A\in\R^{m\times n}\), show that all six matrices are symmetric idempotent:
\[
AA^+,\quad A^+A,\quad I_m-AA^+,\quad I_n-A^+A,
\quad A(A^\top A)^+A^\top,\quad A^\top(AA^\top)^+A.
\]
\end{problem}
\begin{solution}
Write \(A=UDV^\top\) with orthonormal columns in the compact factors.
Then the first two matrices are \(UU^\top,VV^\top\), whose squares
are themselves. Their complements square to themselves because
\((I-P)^2=I-2P+P^2=I-P\), and are symmetric.
Also \((A^\top A)^+=VD^{-2}V^\top\), since the four equations reduce
to the nonzero diagonal entries. Thus
\(A(A^\top A)^+A^\top=UU^\top\). Similarly the sixth matrix is
\(VV^\top\). The zero case gives zero projectors and identity complements.
\end{solution}

\begin{problem}\label{ex:10:self-inverse}
(a) If \(A\) is symmetric idempotent, prove \(A^+=A\).
(b) Does the converse hold?
(c) Show \((AA^+)^+=AA^+\) and \((A^+A)^+=A^+A\).
\end{problem}
\begin{solution}
(a) The choice \(X=A\) gives both triple products \(A^3=A\), and both
products are the symmetric matrix \(A^2=A\); uniqueness proves the claim.
(b) No: for \(A=-I\), the inverse is \(-I=A\), but \(A^2=I\ne A\).
(c) Both matrices are symmetric idempotent by the preceding exercise,
so apply (a) to each.
\end{solution}

\begin{problem}\label{ex:10:mp-ranks}
Show that \(A,A^+,AA^+,A^+A\) have the same rank.
\end{problem}
\begin{solution}
For compact \(A=UDV^\top\) of rank \(r\), the matrix
\(A^+=VD^{-1}U^\top\) has image \(\operatorname{im}V\), of dimension
\(r\). The matrices \(AA^+=UU^\top\) and \(A^+A=VV^\top\) project
onto \(r\)-dimensional spaces, so also have rank \(r\).
For \(A=0\) every rank is zero.
\end{solution}

\begin{problem}\label{ex:10:mp-identities}
Writing \(A^{+\top}=(A^+)^\top\), prove
\begin{enumerate}
\item \(A^\top AA^+=A^\top=A^+AA^\top\);
\item \(A^\top A^{+\top}A^+=A^+=A^+A^{+\top}A^\top\);
\item \((A^\top A)^+=A^+A^{+\top}\);
\item \((AA^\top)^+=A^{+\top}A^+\);
\item \(A^+=(A^\top A)^+A^\top=A^\top(AA^\top)^+\).
\end{enumerate}
\end{problem}
\begin{solution}
Let \(A=UDV^\top\), \(A^+=VD^{-1}U^\top\), with \(U^\top U=V^\top V=I\).
(a) Both products reduce to \(VDU^\top=A^\top\).
(b) Both products reduce to \(VD^{-1}U^\top=A^+\).
(c) \(A^\top A=VD^2V^\top\), whose MP-inverse is \(VD^{-2}V^\top\);
also \(A^+A^{+\top}=VD^{-1}U^\top UD^{-1}V^\top=VD^{-2}V^\top\).
(d) The same multiplication gives \((AA^\top)^+=UD^{-2}U^\top
=A^{+\top}A^+\).
(e) Multiplying the expressions in (c),(d) by \(A^\top=VDU^\top\)
on the stated sides gives \(VD^{-1}U^\top\) in both cases.
For \(A=0\), all displayed products vanish.
\end{solution}

\begin{problem}\label{ex:10:mp-composite}
Prove the following identities, with all products conformable:
\begin{enumerate}
\item \(A(A^\top A)^+A^\top A=A=AA^\top(AA^\top)^+A\).
\item For positive definite \(V\),
\((X^\top V^{-1}X)(X^\top V^{-1}X)^+X^\top=X^\top\).
\item Set \(C=(B^\top B)^+\) and \(K=ABCB^\top A^\top\).
Show \(KK^+AB=AB\).
\end{enumerate}
\end{problem}
\begin{solution}
(a) Each left expression is \(AA^+A=A\), by the identities just proved.
(b) Put \(F=V^{-1/2}X\), \(G=F^\top F=X^\top V^{-1}X\).
Part (a) gives \(FG^+G=F\). Multiply by \(V^{1/2}\) and transpose:
\(GG^+X^\top=X^\top\), since \(G,G^+\) are symmetric.
(c) The matrix \(C\) is positive semidefinite; let \(F=C^{1/2}\) and
\(D=ABF\). Then \(DD^\top=K\), and (a) gives \(KK^+D=D\).
Multiply on the right by \(FB^\top B\). Since
\(B C B^\top B=B\), also by (a), we have
\(DFB^\top B=ABCB^\top B=AB\). Thus \(KK^+AB=AB\), as required.
\end{solution}

\begin{problem}\label{ex:10:full-rank}
Prove \(A^+=(A^\top A)^{-1}A^\top\) for full column rank and
\(A^+=A^\top(AA^\top)^{-1}\) for full row rank.
Deduce \(A^+A=I_n\) in the first case and \(AA^+=I_m\) in the second.
\end{problem}
\begin{solution}
Full column rank gives \(x^\top A^\top Ax=\norm{Ax}^2>0\) for nonzero
\(x\), so \(A^\top A\) is nonsingular. Its MP-inverse is its ordinary
inverse, and the identity \(A^+=(A^\top A)^+A^\top\) gives the formula.
Apply the same argument to \(A^\top\), using
\((A^\top)^+=(A^+)^\top\), for the full row rank case.
Multiplication by \(A\) on the indicated side gives the two identities.
\end{solution}

\begin{problem}\label{ex:10:zero-products}
Prove (a) \(A=0\iff A^+=0\);
(b) \(AB=0\iff B^+A^+=0\);
(c) \(A^+B=0\iff A^\top B=0\).
\end{problem}
\begin{solution}
(a) Zero has zero MP-inverse; the reverse follows from \(A=AA^+A\).
(b) If \(AB=0\), then
\[
B^+A^+=(B^\top B)^+B^\top A^\top(AA^\top)^+=0.
\]
Apply this implication to the pair \(B^+,A^+\) and use
\((A^+)^+=A\), \((B^+)^+=B\) for the reverse.
(c) If \(A^\top B=0\), then \(A^+B=(A^\top A)^+A^\top B=0\).
Conversely \(A^\top B=A^\top AA^+B=0\).
\end{solution}

\begin{problem}\label{ex:10:rank-one}
If \(\operatorname{rank}A=1\), prove
\(A^+=A^\top/\tr(AA^\top)\). In particular determine the MP-inverse
of a column vector, including the zero vector.
\end{problem}
\begin{solution}
Write \(A=ab^\top\) with both vectors nonzero. Set
\(c=\norm{a}^2\norm{b}^2=\tr(AA^\top)>0\), \(X=ba^\top/c\).
Then \(AX=aa^\top/\norm{a}^2\) and \(XA=bb^\top/\norm{b}^2\),
both symmetric. They act as the identity on \(a\), respectively \(b\),
so \(AXA=A\) and \(XAX=X\). Hence \(X=A^+\).
Taking scalar \(b=1\) gives \(a^+=a^\top/(a^\top a)\) for \(a\ne0\).
For a zero \(m\)-column the inverse is a zero \(1\times m\) row.
\end{solution}

\begin{problem}\label{ex:10:block-diagonal}
Prove that the MP-inverse of \(\operatorname{diag}(A_1,A_2)\) is
\(\operatorname{diag}(A_1^+,A_2^+)\), allowing rectangular diagonal blocks.
\end{problem}
\begin{solution}
Call the two displayed block matrices \(A,X\). Block multiplication gives
\(AXA=\operatorname{diag}(A_1A_1^+A_1,A_2A_2^+A_2)=A\)
and \(XAX=X\). The products \(AX\) and \(XA\) are block diagonal
with symmetric blocks \(A_iA_i^+\) and \(A_i^+A_i\), respectively.
Thus all four equations hold, and uniqueness identifies \(X=A^+\).
\end{solution}

\begin{problem}\label{ex:10:symmetric-inverse}
Let \(A=S\Lambda S^\top\) be symmetric of rank \(r>0\), where
\(S^\top S=I_r\) and \(\Lambda\) is diagonal with nonzero entries,
possibly negative. Prove \(A^+=S\Lambda^{-1}S^\top\).
\end{problem}
\begin{solution}
Put \(X=S\Lambda^{-1}S^\top\). Both products \(AX,XA\) equal
\(SS^\top\), a symmetric matrix. Also
\(AXA=SS^\top S\Lambda S^\top=A\) and
\(XAX=SS^\top S\Lambda^{-1}S^\top=X\).
Thus \(X\) satisfies all four equations, regardless of the signs of the
nonzero diagonal entries.
\end{solution}

\needspace{8\baselineskip}
\subsection*{Products, projections, and ranges}

\begin{problem}\label{ex:10:reverse-order}
Prove the following reverse-order formulas:
\begin{enumerate}
\item If \(BB^\top=I\), then \((AB)^+=B^\top A^+\).
\item If \(A\) has full column rank and \(B\) full row rank, then
\((AB)^+=B^+A^+\).
\item If \(A=BC^\top\), with \(B,C\) both having \(r>0\) independent
columns, then
\[A^+=C(B^\top BC^\top C)^{-1}B^\top.\]
\end{enumerate}
\end{problem}
\begin{solution}
(a) Put \(X=B^\top A^+\). The triple products reduce to
\(ABXAB=AA^+AB=AB\) and \(XABX=B^\top A^+AA^+=X\).
The two products are \(ABX=AA^+\) and \(XAB=B^\top A^+AB\),
both symmetric; thus \(X=(AB)^+\).
(b) Put \(X=B^+A^+\). Since \(A^+A=I\), \(BB^+=I\), the triple
products give \(ABXAB=AB\), \(XABX=X\). The two products reduce to
\(AA^+\) and \(B^+B\), both symmetric.
(c) Apply (b) to \(B,C^\top\) and the full-rank formulas:
\[
A^+=C(C^\top C)^{-1}(B^\top B)^{-1}B^\top
=C(B^\top BC^\top C)^{-1}B^\top.
\]
Both Gram matrices are positive definite, so their inverses exist.
\end{solution}

\begin{problem}\label{ex:10:equal-projectors}
Let \(A\) be square. (a) If \(A\) is symmetric, show \(AA^+=A^+A\).
(b) Does the converse hold? (c) Prove that equality holds exactly when
\(\operatorname{im}A=\operatorname{im}A^\top\).
\end{problem}
\begin{solution}
(a) In a compact spectral decomposition \(A=S\Lambda S^\top\), both
products are \(SS^\top\); the zero case is immediate.
(b) No: \(\begin{pmatrix}1&1\\0&1\end{pmatrix}\) is nonsymmetric
and nonsingular, so both products are \(I_2\).
(c) The two products are the orthogonal projectors onto the two indicated
images. Equality of matrices implies equality of their images.
Conversely, if the images agree, uniqueness of the orthogonal projection
onto that common subspace gives equality of the matrices.
\end{solution}

\begin{problem}\label{ex:10:psd-factor}
If \(A=BB^\top\), where \(B\) has \(r>0\) independent columns, prove
\(A^+=B(B^\top B)^{-2}B^\top\).
\end{problem}
\begin{solution}
The full-rank reverse-order formula for \(BB^\top\) gives
\[
A^+=(B^\top)^+B^+
=B(B^\top B)^{-1}(B^\top B)^{-1}B^\top.
\]
Here \(B^\top B\) is positive definite, so the stated expression exists.
For a rank-zero positive semidefinite matrix, the separate answer is \(A^+=0\).
\end{solution}

\begin{problem}\label{ex:10:normal-equivalence}
Prove \(A^\top AB=A^\top C\iff AB=AA^+C\).
\end{problem}
\begin{solution}
If the first equality holds, premultiply by \(A(A^\top A)^+\):
\[
AB=A(A^\top A)^+A^\top AB=A(A^\top A)^+A^\top C=AA^+C.
\]
The two simplifications use the preceding MP identities.
Conversely multiply \(AB=AA^+C\) by \(A^\top\) and use
\(A^\top AA^+=A^\top\). This gives \(A^\top AB=A^\top C\).
\end{solution}

\begin{problem}\label{ex:10:range-product}
If \(B\) has full row rank, prove \((AB)(AB)^+=AA^+\).
Is full row rank necessary for this equality?
\end{problem}
\begin{solution}
The inclusion \(\operatorname{im}(AB)\subseteq\operatorname{im}A\)
always holds. Full row rank gives \(BB^+=I\), hence \(A=(AB)B^+\),
which gives the reverse inclusion. The two orthogonal projectors
therefore agree. The hypothesis is not necessary: let
\(A=B=\operatorname{diag}(1,0)\). Then \(AB=A\) and both projectors
equal \(A\), although \(B\) has rank one rather than two.
\end{solution}

\begin{problem}\label{ex:10:nested-projectors}
Let \(A\) be symmetric idempotent. (a) If \(AB=B\), show that
\(A-BB^+\) is symmetric idempotent of rank
\(\operatorname{rank}A-\operatorname{rank}B\).
(b) Does the converse hold?
\end{problem}
\begin{solution}
Put \(P=BB^+\). From \(AB=B\), \(AP=P\), and transposition gives
\(PA=P\). Therefore \((A-P)^2=A-AP-PA+P=A-P\), and symmetry is immediate.
Its image is \(\operatorname{im}A\cap(\operatorname{im}B)^\perp\):
\(A(A-P)=A-P\), \(P(A-P)=0\); conversely \(Az=z,Pz=0\) gives
\((A-P)z=z\). Since \(\operatorname{im}B\subseteq\operatorname{im}A\),
orthogonal decomposition within \(\operatorname{im}A\) gives the rank difference.
(b) Yes; idempotence already suffices. If \(C=A-P\) is idempotent, it is
symmetric, so \(z^\top Cz=\norm{Cz}^2\geqslant0\).
For \(z\in\operatorname{im}P\), this says
\(z^\top Az\geqslant\norm{z}^2\). But \(A\) is an orthogonal projector,
so \(z^\top Az=\norm{Az}^2\leqslant\norm{z}^2\).
Equality and Pythagoras imply \((I-A)z=0\). Thus \(A\) fixes every
column of \(B\), giving \(AB=B\).
\end{solution}

\begin{problem}\label{ex:10:identify-projector}
For a symmetric idempotent \(n\times n\) matrix \(A\), show:
(a) if \(AB=B\) and \(\operatorname{rank}A=\operatorname{rank}B\), then
\(A=BB^+\);
(b) if \(AB=0\) and \(\operatorname{rank}A+\operatorname{rank}B=n\),
then \(A=I_n-BB^+\).
\end{problem}
\begin{solution}
(a) The preceding exercise makes \(A-BB^+\) a matrix of rank zero,
so it vanishes.
(b) The matrix \(C=I-A\) is symmetric idempotent and fixes \(B\).
Its image is \(\ker A\), so its rank is
\(n-\operatorname{rank}A=\operatorname{rank}B\). Part (a) applied to
\(C\) gives \(I-A=BB^+\).
\end{solution}

\needspace{8\baselineskip}
\subsection*{Vector equations and least squares}

\begin{problem}\label{ex:10:homogeneous}
(a) Show that every solution of \(Ax=0\) has the form
\(x=(I-A^+A)q\), with \(q\) arbitrary, and that every such vector is
a solution. (b) Determine when the solution is unique.
(c) Prove that otherwise there are infinitely many solutions.
\end{problem}
\begin{solution}
Multiplication gives \(A(I-A^+A)q=(A-AA^+A)q=0\).
Conversely if \(Ax=0\), then \((I-A^+A)x=x\), so choose \(q=x\).
The solution is unique exactly when \(\ker A=\{0\}\), exactly when
\(A\) has full column rank; it is then zero. Otherwise choose nonzero
\(z\in\ker A\). The vectors \(tz\), \(t\in\R\), are distinct solutions,
proving infinitude.
\end{solution}

\begin{problem}\label{ex:10:inconsistent}
Give a nonhomogeneous real system \(Ax=b\), \(b\ne0\), with no solution.
\end{problem}
\begin{solution}
Take \(A=\begin{pmatrix}1&2\\1&2\end{pmatrix}\),
\(b=(1,2)^\top\). The two equations would require the same quantity
\(x_1+2x_2\) to equal both one and two, a contradiction.
\end{solution}

\begin{problem}\label{ex:10:consistency}
(a) Show that the following are equivalent: \(Ax=b\) is consistent;
\(b\in\operatorname{im}A\); \(\operatorname{rank}(A\mid b)=\operatorname{rank}A\);
and \(AA^+b=b\).

(b) Does full column rank of \(A\) guarantee consistency of \(Ax=b\)?
Give an explicit counterexample if it does not.
\end{problem}
\begin{solution}
(a) The first two are equivalent because \(Ax\) runs through the column
space of \(A\). Appending \(b\) leaves the column-space dimension
unchanged exactly when \(b\) was already in that space; otherwise it
increases the dimension by one. This proves the third equivalence.
Finally \(AA^+\) projects onto that space, and fixes exactly its members.
Equivalently, consistency gives \(b=Ax=AA^+Ax=AA^+b\), while
\(AA^+b=b\) supplies the solution \(x=A^+b\).

(b) No. Take \(A=(1,0)^\top\) and \(b=(0,1)^\top\).
The single column of \(A\) is nonzero, so \(A\) has full column rank.
But \(Ax=(x,0)^\top\) cannot equal \(b\), whose second coordinate is one.
\end{solution}

\begin{problem}\label{ex:10:all-exact}
If \(Ax=b\) is consistent, prove that its complete solution set is
\(x=A^+b+(I-A^+A)q\), where \(q\) is arbitrary.
\end{problem}
\begin{solution}
Consistency gives \(AA^+b=b\), so \(x_0=A^+b\) is one solution.
Then \(Ax=b\iff A(x-x_0)=0\). The complete kernel formula in
Exercise~\ref{ex:10:homogeneous} gives exactly
\(x-x_0=(I-A^+A)q\). The same calculation verifies every proposed vector.
\end{solution}

\begin{problem}\label{ex:10:ls-identity}
Whether or not \(Ax=b\) is consistent, prove
\[
\norm{Ax-b}^2=(x-A^+b)^\top A^\top A(x-A^+b)
+b^\top(I-AA^+)b.
\]
Deduce that \(A^+b\) minimizes the residual norm.
\end{problem}
\begin{solution}
Let \(P=AA^+\). Write
\(Ax-b=A(x-A^+b)-(I-P)b\).
The two terms are perpendicular, since \(A^\top(I-P)=0\).
Their squared norms add. Because \(I-P\) is symmetric idempotent,
\(\norm{(I-P)b}^2=b^\top(I-P)b\), proving the formula.
The first term is \(\norm{A(x-A^+b)}^2\geqslant0\), and equals zero
at \(x=A^+b\). Thus this choice attains the global minimum.
\end{solution}

\begin{problem}\label{ex:10:surjective-injective}
Prove (a) that \(Ax=b\) is consistent for every \(b\in\R^m\) exactly
when \(A\) has full row rank; (b) that a consistent \(Ax=b\) has a
unique solution exactly when \(A\) has full column rank.
\end{problem}
\begin{solution}
(a) Consistency for every \(b\) means \(\operatorname{im}A=\R^m\),
equivalent to rank \(m\), or full row rank.
(b) Fix one solution \(x_0\). All solutions are \(x_0+\ker A\), so
uniqueness is equivalent to \(\ker A=\{0\}\). Rank–nullity makes this
equivalent to rank \(n\), or full column rank. Neither conclusion alone
requires \(A\) to be square.
\end{solution}

\needspace{8\baselineskip}
\subsection*{Matrix equations}

\begin{problem}\label{ex:10:matrix-all}
Let \(A\in\R^{m\times n}\), \(B\in\R^{p\times q}\), and
\(C\in\R^{m\times q}\). Prove that \(AXB=C\) is consistent exactly
when \(AA^+CB^+B=C\). If consistent, prove that all solutions are
\[
X=A^+CB^++Q-A^+AQBB^+,\qquad Q\in\R^{n\times p}.
\]
\end{problem}
\begin{solution}
If \(AXB=C\), multiplication gives
\(AA^+CB^+B=AA^+AXB B^+B=AXB=C\).
Conversely that equality shows that \(X_0=A^+CB^+\) is a solution.
For arbitrary \(Q\),
\(A(Q-A^+AQBB^+)B=AQB-AA^+AQBB^+B=0\),
so every displayed matrix is a solution.
If \(X\) is any solution, set \(Y=X-X_0\), so \(AYB=0\).
Then \(A^+AYBB^+=0\), and choosing \(Q=Y\) in the formula returns
\(X_0+Y=X\). Thus no solutions are missing.
\end{solution}

\begin{problem}\label{ex:10:matrix-uniqueness}
For the same positive matrix dimensions, prove:
(a) \(AXB=C\) is consistent for every \(C\) exactly when \(A\) has
full row rank and \(B\) full column rank;
(b) when consistent, its solution is unique exactly when \(A\) has
full column rank and \(B\) full row rank.
\end{problem}
\begin{solution}
(a) The two full-rank conditions give \(AA^+=I_m\), \(B^+B=I_q\),
so the preceding consistency test holds for every \(C\).
Conversely, if \(A\) lacks full row rank, choose nonzero \(u\in\ker A^\top\).
Then \(u^\top AXB=0\), whereas \(C=ue_1^\top\) has
\(u^\top C=\norm{u}^2e_1^\top\ne0\). If \(B\) lacks full column rank,
choose nonzero \(v\in\ker B\). Every \(AXB\) annihilates \(v\), but
\(C=e_1v^\top\) does not.
(b) If both stated ranks are full, multiplying a homogeneous equation
\(AYB=0\) by \(A^+\) and \(B^+\) gives \(Y=0\), proving uniqueness.
If \(A\) has a nonzero kernel vector \(z\), then \(Y=ze_1^\top\ne0\)
satisfies \(AYB=0\). If \(B\) lacks full row rank, choose nonzero
\(w\in\ker B^\top\); then \(Y=e_1w^\top\ne0\) also satisfies it.
Adding either such \(Y\) to a solution disproves uniqueness.
\end{solution}

\begin{problem}\label{ex:10:matrix-homogeneous}
Find all matrices solving (a) \(AX=0\); (b) \(XA=0\).
\end{problem}
\begin{solution}
(a) All solutions are \(X=(I-A^+A)Q\), with arbitrary conformable \(Q\).
Multiplication by \(A\) gives zero. Conversely \(AX=0\) implies
\((I-A^+A)X=X\), so choose \(Q=X\).
(b) All solutions are \(X=Q(I-AA^+)\). Each annihilates \(A\), since
\((I-AA^+)A=0\). Conversely \(XA=0\) gives \(X(I-AA^+)=X\).
\end{solution}

\needspace{8\baselineskip}
\subsection*{Generalized inverses}

\begin{problem}\label{ex:10:generalized-all}
A generalized inverse of \(A\) is any matrix \(G\) satisfying \(AGA=A\).
Prove existence and show that all such matrices are
\[
G=A^++Q-A^+AQAA^+,\qquad Q\text{ arbitrary and conformable}.
\]
\end{problem}
\begin{solution}
The choice \(G=A^+\) establishes existence. For the displayed formula,
\[
AGA=AA^+A+AQA-AA^+AQAA^+A=A+AQA-AQA=A.
\]
Conversely, if \(AGA=A\), set \(Q=G-A^+\). Then \(AQA=0\), so
\(A^+AQAA^+=0\), and the displayed expression becomes \(A^++Q=G\).
Thus it gives precisely all generalized inverses.
\end{solution}

\begin{problem}\label{ex:10:generalized-projector}
For any generalized inverse \(G\) of \(A\), prove that \(GA\) is
idempotent. Show that it need not be symmetric. If it is symmetric,
prove that \(GA=A^+A\), so this product is unique.
\end{problem}
\begin{solution}
The identity \(AGA=A\) gives \((GA)^2=G(AGA)=GA\).
For a counterexample, take
\(A=\operatorname{diag}(1,0)\),
\(G=\begin{pmatrix}1&0\\1&0\end{pmatrix}\).
Then \(AGA=A\), but \(GA=G\) is nonsymmetric.
If \(GA\) is symmetric, it is an orthogonal projector.
Also \(\ker(GA)=\ker A\): one inclusion is immediate; if \(GAx=0\),
then \(Ax=AGAx=0\). Thus it projects onto \((\ker A)^\perp\), exactly
the image of the orthogonal projector \(A^+A\). Uniqueness gives equality.
\end{solution}

\begin{problem}\label{ex:10:generalized-fit}
Let \(G\) be any generalized inverse of \(A^\top A\).
Prove \(AGA^\top=A(A^\top A)^+A^\top\).
Conclude that this matrix is unique, symmetric, and idempotent.
\end{problem}
\begin{solution}
For nonzero \(A=UDV^\top\), the generalized-inverse equation is
\(VD^2V^\top GVD^2V^\top=VD^2V^\top\).
Multiply by \(V^\top\) on the left and \(V\) on the right to get
\(D^2(V^\top GV)D^2=D^2\), hence \(V^\top GV=D^{-2}\).
Therefore
\[
AGA^\top=UD(V^\top GV)DU^\top=UU^\top
=A(A^\top A)^+A^\top.
\]
This orthogonal projector is independent of \(G\), symmetric, and idempotent.
If \(A=0\), the expression is zero for every \(G\), with the same conclusions.
\end{solution}

\begin{problem}\label{ex:10:generalized-rank}
Let \(A\) be symmetric idempotent.
(a) Show that all its generalized inverses are \(A+Q-AQA\).
(b) Show that \(B=A+\alpha(I-A)\) is a generalized inverse for every \(\alpha\).
(c) Prove that \(B\) is nonsingular for \(\alpha\ne0\).
(d) Conclude that a generalized inverse need not have the same rank as \(A\).
\end{problem}
\begin{solution}
(a) In the general formula for generalized inverses, \(A^+=A\) and
\(A^2=A\), giving \(A+Q-AQA\).
(b) Take \(Q=\alpha I\); the formula gives the stated \(B\).
(c) Since \(A(I-A)=(I-A)A=0\), multiplication on either side gives
\[
[A+\alpha(I-A)][A+\alpha^{-1}(I-A)]=A+(I-A)=I.
\]
The factors commute, so the second is a two-sided inverse.
(d) If \(A\) is singular and \(\alpha\ne0\), then \(B\) has full rank
while \(A\) does not. For example \(A=\operatorname{diag}(1,0)\),
\(\alpha=1\) gives \(B=I_2\).
\end{solution}

\begin{problem}\label{ex:10:generalized-solutions}
For any generalized inverse \(G\) of \(A\), prove that \(Ax=b\) is
consistent exactly when \(AGb=b\). When consistent, show that all
solutions are \(x=Gb+(I-GA)q\), \(q\) arbitrary.
\end{problem}
\begin{solution}
If \(b=Ax\), then \(AGb=AGAx=Ax=b\). Conversely \(AGb=b\) exhibits
\(Gb\) as a solution. Under consistency,
\[
A[Gb+(I-GA)q]=b+(A-AGA)q=b,
\]
so every displayed vector is a solution. If \(x\) is any solution,
choose \(q=x\): \(Gb+(I-GA)x=Gb+x-Gb=x\).
Thus the parametrization is complete. No symmetry of \(AG\) or \(GA\)
is required for these exact-system statements.
\end{solution}
